\section{Método de lectura de código: primero ver cómo fluyen los datos hacia el Trainer}
El valor de esta clase está en fijar con claridad las dos abstracciones centrales, Dataset y Trainer. Al leer código de un framework de verdad, lo más útil de seguir primero es: cómo pasa una muestra de \texttt{get\_item} a un batch, y de ahí fluye hacia el Trainer para completar la propagación hacia adelante y hacia atrás.

\begin{knowledgebox}{La ruta más corta para leer un framework de entrenamiento}
\begin{itemize}
  \item Primero ver cómo organiza el Dataset las entradas y salidas
  \item Después ver cómo consume el Trainer el batch
  \item Al final ver cómo envuelve la estrategia distribuida el ciclo de entrenamiento
\end{itemize}
\end{knowledgebox}

\subsection{Resumen de la sección}
Para entender un framework de entrenamiento no es necesario empezar por el código distribuido más complejo; ver primero el flujo de datos y la lógica del Trainer suele facilitar más la construcción de una visión de conjunto.

\newpage
